\documentclass[12pt,a4paper]{article}
\usepackage[utf8]{inputenc}
\usepackage[T1]{fontenc}
\usepackage[french]{babel}
\usepackage{amsmath, amssymb}
\usepackage{graphicx}
\usepackage{geometry}
\geometry{margin=2cm}
\usepackage{enumitem}
\usepackage{float}
\usepackage{fancyhdr}
\usepackage{titlesec}
\pagestyle{fancy}

% Nettoyer les styles par défaut
\fancyhf{}

% En-têtes
\fancyhead[L]{Fiches d'exercices pour Maximilien (2)} % à gauche
\fancyhead[R]{Matteo Dubresson} % à droite

% Pieds de page
\fancyfoot[C]{\thepage} % numéro de page centré

%Titre et généralités
\title{Fiche d'exercices \\ pour Maximilien (2)}
\date{19/09/2025}
\author{}

%Espace entre les subsections
\titlespacing*{\subsection}{0pt}{5ex plus 0.5ex minus 0.2ex}{1.5ex plus 0.2ex}




\begin{document}

\maketitle

\subsection*{Exercice 1 :}

Soit $(u_n)_{n \in \mathbb{N}}$ la suite définie par $u_0 = 2$ et pour tout $n \in \mathbb{N}$ par :
\[
u_{n+1} = 5u_n - 4
\]
Montrer que pour tout $n \in \mathbb{N}$ :
\[
u_n = 5^n + 1
\]

\bigskip

\subsection*{Exercice 2 :}

Soit $(y_n)$ la suite définie par $y_0 = 1$, $y_1 = 2$ et pour tout $n \in \mathbb{N}$ :
\[
y_{n+2} = 3y_{n+1} - 2y_n
\]

\begin{enumerate}
    \item Calculer $y_2$ et $y_3$.
    \item Montrer par récurrence forte que pour tout $n \in \mathbb{N}$ :
    \[
    y_n = 2^n
    \]
\end{enumerate}

\bigskip


\bigskip

\subsection*{Exercice 3 :}

Soient $P$ et $Q$ deux propositions logiques.

\begin{enumerate}
    \item Dresser la table de vérité de la proposition $(P \implies Q) \iff (\neg P \lor Q)$.
    \item En déduire la négation formelle de l'implication $P \implies Q$.
\end{enumerate}

\bigskip

\subsection*{Exercice 5 :}

Soit $f : \mathbb{R} \to \mathbb{R}$ une fonction. Écrire la négation des propositions suivantes :

\begin{enumerate}
    \item $\forall x \in \mathbb{R}, \quad f(x) > 0$
    \item $\exists M \in \mathbb{R}, \quad \forall x \in \mathbb{R}, \quad |f(x)| \le M$
    \item $\forall x \in \mathbb{R}, \quad \big( f(x) = 0 \implies x \ge 0 \big)$
\end{enumerate}

\bigskip

\subsection*{Exercice 6 :}

Soit $n \in \mathbb{Z}$. Montrer par contraposée l'assertion suivante :
\[
n^2 \text{ est pair} \implies n \text{ est pair}
\]


\subsection*{Exercice 7 :}

Soit $P$ une proposition logique.
\begin{enumerate}
    \item Écrire la table de vérité de la proposition $(P \land \neg P)$.
    \item Comment appelle-t-on ce type de proposition en logique ?
\end{enumerate}

\bigskip

\subsection*{Exercice 8 :}

Soient $P$, $Q$ et $R$ trois propositions logiques.
\begin{enumerate}
    \item Montrer, à l'aide d'une table de vérité ou par le calcul algébrique, l'équivalence suivante (loi de distributivité) :
    \[
    P \land (Q \lor R) \iff (P \land Q) \lor (P \land R)
    \]
    \item Énoncer la loi de Morgan reliant $\neg (P \land Q)$ aux propositions $\neg P$ et $\neg Q$.
\end{enumerate}

\bigskip

\subsection*{Exercice 9 :}

Soit $n \in \mathbb{N}$. On considère la proposition $A(n) : « \, n^2 - 1 \text{ est divisible par } 3 \, »$.
\begin{enumerate}
    \item La proposition $A(0)$ est-elle vraie ?
    \item Exprimer à l'aide des quantificateurs la proposition : « Pour tout entier naturel $n$, $A(n)$ est vraie ».
    \item Donner la négation de cette proposition quantifiée, puis déterminer si la proposition initiale est vraie ou fausse en justifiant par un contre-exemple.
\end{enumerate}


\end{document}